% Options for packages loaded elsewhere
% Options for packages loaded elsewhere
\PassOptionsToPackage{unicode}{hyperref}
\PassOptionsToPackage{hyphens}{url}
%
\documentclass[
  ignorenonframetext,
]{beamer}
\newif\ifbibliography
\usepackage{pgfpages}
\setbeamertemplate{caption}[numbered]
\setbeamertemplate{caption label separator}{: }
\setbeamercolor{caption name}{fg=normal text.fg}
\beamertemplatenavigationsymbolsempty
% remove section numbering
\setbeamertemplate{part page}{
  \centering
  \begin{beamercolorbox}[sep=16pt,center]{part title}
    \usebeamerfont{part title}\insertpart\par
  \end{beamercolorbox}
}
\setbeamertemplate{section page}{
  \centering
  \begin{beamercolorbox}[sep=12pt,center]{section title}
    \usebeamerfont{section title}\insertsection\par
  \end{beamercolorbox}
}
\setbeamertemplate{subsection page}{
  \centering
  \begin{beamercolorbox}[sep=8pt,center]{subsection title}
    \usebeamerfont{subsection title}\insertsubsection\par
  \end{beamercolorbox}
}
% Prevent slide breaks in the middle of a paragraph
\widowpenalties 1 10000
\raggedbottom
\AtBeginPart{
  \frame{\partpage}
}
\AtBeginSection{
  \ifbibliography
  \else
    \frame{\sectionpage}
  \fi
}
\AtBeginSubsection{
  \frame{\subsectionpage}
}
\usepackage{iftex}
\ifPDFTeX
  \usepackage[T1]{fontenc}
  \usepackage[utf8]{inputenc}
  \usepackage{textcomp} % provide euro and other symbols
\else % if luatex or xetex
  \usepackage{unicode-math} % this also loads fontspec
  \defaultfontfeatures{Scale=MatchLowercase}
  \defaultfontfeatures[\rmfamily]{Ligatures=TeX,Scale=1}
\fi
\usepackage{lmodern}

\usetheme[]{metropolis}
\ifPDFTeX\else
  % xetex/luatex font selection
\fi
% Use upquote if available, for straight quotes in verbatim environments
\IfFileExists{upquote.sty}{\usepackage{upquote}}{}
\IfFileExists{microtype.sty}{% use microtype if available
  \usepackage[]{microtype}
  \UseMicrotypeSet[protrusion]{basicmath} % disable protrusion for tt fonts
}{}
\makeatletter
\@ifundefined{KOMAClassName}{% if non-KOMA class
  \IfFileExists{parskip.sty}{%
    \usepackage{parskip}
  }{% else
    \setlength{\parindent}{0pt}
    \setlength{\parskip}{6pt plus 2pt minus 1pt}}
}{% if KOMA class
  \KOMAoptions{parskip=half}}
\makeatother


\usepackage{longtable,booktabs,array}
\usepackage{calc} % for calculating minipage widths
\usepackage{caption}
% Make caption package work with longtable
\makeatletter
\def\fnum@table{\tablename~\thetable}
\makeatother
\usepackage{graphicx}
\makeatletter
\newsavebox\pandoc@box
\newcommand*\pandocbounded[1]{% scales image to fit in text height/width
  \sbox\pandoc@box{#1}%
  \Gscale@div\@tempa{\textheight}{\dimexpr\ht\pandoc@box+\dp\pandoc@box\relax}%
  \Gscale@div\@tempb{\linewidth}{\wd\pandoc@box}%
  \ifdim\@tempb\p@<\@tempa\p@\let\@tempa\@tempb\fi% select the smaller of both
  \ifdim\@tempa\p@<\p@\scalebox{\@tempa}{\usebox\pandoc@box}%
  \else\usebox{\pandoc@box}%
  \fi%
}
% Set default figure placement to htbp
\def\fps@figure{htbp}
\makeatother

\ifLuaTeX
  \usepackage{luacolor}
  \usepackage[soul]{lua-ul}
\else
  \usepackage{soul}
  \makeatletter
  \let\HL\hl
  \renewcommand\hl{% fix for beamer highlighting
    \let\set@color\beamerorig@set@color
    \let\reset@color\beamerorig@reset@color
    \HL}
  \makeatother
\fi




\setlength{\emergencystretch}{3em} % prevent overfull lines

\providecommand{\tightlist}{%
  \setlength{\itemsep}{0pt}\setlength{\parskip}{0pt}}



 


\setbeamerfont{title}{size=\large} \usepackage{hyperref} \setbeamertemplate{footline}[frame number]
\makeatletter
\@ifpackageloaded{caption}{}{\usepackage{caption}}
\AtBeginDocument{%
\ifdefined\contentsname
  \renewcommand*\contentsname{Table of contents}
\else
  \newcommand\contentsname{Table of contents}
\fi
\ifdefined\listfigurename
  \renewcommand*\listfigurename{List of Figures}
\else
  \newcommand\listfigurename{List of Figures}
\fi
\ifdefined\listtablename
  \renewcommand*\listtablename{List of Tables}
\else
  \newcommand\listtablename{List of Tables}
\fi
\ifdefined\figurename
  \renewcommand*\figurename{Figure}
\else
  \newcommand\figurename{Figure}
\fi
\ifdefined\tablename
  \renewcommand*\tablename{Table}
\else
  \newcommand\tablename{Table}
\fi
}
\@ifpackageloaded{float}{}{\usepackage{float}}
\floatstyle{ruled}
\@ifundefined{c@chapter}{\newfloat{codelisting}{h}{lop}}{\newfloat{codelisting}{h}{lop}[chapter]}
\floatname{codelisting}{Listing}
\newcommand*\listoflistings{\listof{codelisting}{List of Listings}}
\makeatother
\makeatletter
\makeatother
\makeatletter
\@ifpackageloaded{caption}{}{\usepackage{caption}}
\@ifpackageloaded{subcaption}{}{\usepackage{subcaption}}
\makeatother

\usepackage{bookmark}
\IfFileExists{xurl.sty}{\usepackage{xurl}}{} % add URL line breaks if available
\urlstyle{same}
\hypersetup{
  pdfauthor={Giorgio Arcara},
  hidelinks,
  pdfcreator={LaTeX via pandoc}}


\author{Giorgio Arcara}
\date{}

\begin{document}


\begin{frame}
% TITLE SLIDES
\title{Metodologia della Ricerca Clinica\\ \vspace{1em} \emph{05 - Valutare deficit e danni - Teoria}}
\author{Giorgio Arcara,\\ Università di Padova \\ IRCCS San Camillo, Venezia}

\titlegraphic{

\vspace*{7cm}
\includegraphics[scale=0.15]{Figures/LogoSanCamilloIRCCS_Unipd_alpha.png}
\begin{center}
\vspace{-2.2em}
\includegraphics[scale=0.15]{Figures/CC_license_3_0.png}
\end{center}
}
%\date{\today}
\maketitle
\end{frame}

\section{5. Valutare i deficit e danni
cognitivi}\label{valutare-i-deficit-e-danni-cognitivi}

\begin{frame}{Misurazioni 1/2}
\phantomsection\label{misurazioni-12}
\center

\includegraphics[width=0.8\linewidth,height=\textheight,keepaspectratio]{Figures/thermometer_1.png}
\end{frame}

\begin{frame}{Misurazioni 2/2}
\phantomsection\label{misurazioni-22}
\center

\includegraphics[width=0.8\linewidth,height=\textheight,keepaspectratio]{Figures/thermometer_2.png}
\end{frame}

\begin{frame}{Valutare i deficit cognitivi}
\phantomsection\label{valutare-i-deficit-cognitivi}
Nelle prossime slide userò spesso i termini deficit e danni. \textbf{Non
sono sinonimi}:

\textbf{Deficit (cognitivo)}: una mancanza/difetto funzionale sul piano
cognitivo quando paragonato con un gruppo di riferimento (spesso i pari)
o una traiettoria di sviluppo attesa. Es. una persona che non è in grado
di leggere (con vista integra) ha un deficit cognitivo.

\pause

\textbf{Danno (cognitivo)}: un peggioramento di un'abilità cognitiva
rispetto ad una condizione precedentemente raggiunta, in seguito ad un
evento (può essere puntuale nel tempo, come un ictus, o non
localizzabile, es. degenerazione rispetto ad Alzheimer). \vspace{2em}

\footnotesize

NOTA: questa non è necessariamente una terminologia condivisa. L'aspetto
importante è però distinguere concettualmente questi due tipi di
possibilità
\end{frame}

\begin{frame}{Danno in contesto forense}
\phantomsection\label{danno-in-contesto-forense}
In ambito forense (Italiano) il danno psichico è tradizionalmente
considerato come una sottospecie del danno biologico, ossia della
lesione dell'integrità psicofisica, medicalmente accertabile.

Nel danno in ambito forense, è fondamentale identificare non solo che
sia avvenuto un peggioramento, \textbf{ma che il danno sia anche
\emph{causato} dal fatto esaminato.}
\end{frame}

\begin{frame}{Significato dei punteggi}
\phantomsection\label{significato-dei-punteggi}
Supponiamo di somministrare il test di Linguaggio Figurato 1 di APACS
(perché siamo soddisfatti di sue Validità e Affidabilità) e supponiamo
di ottenere il punteggio 10/15 (quindi 75\% di risposte corrette). Cosa
ci dice questo punteggio? \pause

\begin{center}
  \textbf{I punteggi ai test in sé sono poco informativi}
\end{center}
\pause
\end{frame}

\begin{frame}{Significato dei punteggi 1/2}
\phantomsection\label{significato-dei-punteggi-12}
\begin{center}
\begin{figure}
\includegraphics[scale=0.06]{Figures/triangle.png}
\end{figure}
\end{center}

I punteggi grezzi sono poco informativi non solo perché i valori
dipendono dal range dei punteggi (che spesso è arbitrario), ma anche
perché un punteggio alto o basso dipende molto da quanto il test è
difficile (o facile).

\textbf{Per trarre informazioni utili, in neuropsicologia i punteggi si confrontano spesso con dei dati di riferimento che ci aiutano a classificare la prestazione}.
\end{frame}

\begin{frame}{Trasformazione punteggi e cut-off 2/2}
\phantomsection\label{trasformazione-punteggi-e-cut-off-22}
\begin{center}
\begin{figure}
\includegraphics[scale=0.06]{Figures/triangle.png}
\end{figure}
\end{center}

Per trasformare i punteggi ad un test cognitivo in qualcosa di
interpretabile e utile spesso si effettua una trasformazione dal
punteggio grezzo \underline{alla percentuale di popolazione}
\underline{che ottiene quel punteggio}.

Tra i punteggi trasformati dati spesso rivestono particolare rilevanza
delle soglie specifiche, chiamate \textbf{cut-off}, che nella versione
più diffusa, delimita il 5\% delle prestazioni peggiori.
\end{frame}

\begin{frame}{Cut-off e statistica}
\phantomsection\label{cut-off-e-statistica}
Anche se spesso non è chiaro, il confronto di un punteggio con un
cut-off si configura come l'applicazione di un \textbf{test statistico}
per arrivare ad una decisione (inferenza satistica).

I test test statistici utilizzati in contesto clinico sono spesso non
convenzionali perché \emph{applicati al caso singolo} (mentre gran parte
della statistica inferenziale stima effetti medi riferiti a gruppi).
\end{frame}

\begin{frame}{Dati normativi e cut-off}
\phantomsection\label{dati-normativi-e-cut-off}
I cut-off sono un punteggi-soglia al di sotto del quale la prestazione è
considerata deficitaria e (a seconda del contesto) il riflesso di un
danno cognitivo.

I cut-off vengono calcolati a partire dai dati normativi e spesso
corrisponde al \textbf{5\% dei punteggi peggiori}

Spesso i cut-off solitamente riportati in tabelle, indipendentemente dal
metodo usato per calcolarli. Di recente sono sempre più diffusi software
che trasformano i punteggi e confrontano in maniera automatica con i
cut-offs.
\end{frame}

\begin{frame}{Esempio di cut-off riportati in tabelle}
\phantomsection\label{esempio-di-cut-off-riportati-in-tabelle}
\begin{center}
\includegraphics[width=0.9\linewidth,height=\textheight,keepaspectratio]{Figures/cutoff_humor.png}
\end{center}

\begin{flushright}
  Da APACS (Arcara \& Bambini, 2016)
\end{flushright}
\end{frame}

\begin{frame}{Valutare i deficit/danni cognitivi}
\phantomsection\label{valutare-i-deficitdanni-cognitivi}
\begin{center}
  \textbf{I  cut-off e i dati normativi}
\end{center}

I dati normativi sono dati raccolti su soggetti che assumiamo siano
normali. Servono a stabilire quali prestazioni possiamo considerare
normali e quali no. \vspace{2em}

Chiamiamo i cut-off calcolato sui dati normativi \textbf{cut-off di
normalità}.
\end{frame}

\begin{frame}{Valutare i deficit/danni cognitivi}
\phantomsection\label{valutare-i-deficitdanni-cognitivi-1}
\begin{center}
  \textbf{I  cut-off e i dati normativi}
\end{center}

Sebbene cut-off e dati normativi siano utilizzati quotidianamente in
neuropsicologia clinica, spesso c'è una conoscenza superficiale e
confusa del loro significato. \vspace{2em}

Questa superficialità si ripercuote nelle conclusioni tratte dai
risultati ai test.
\end{frame}

\begin{frame}{Valutare i deficit cognitivi}
\phantomsection\label{valutare-i-deficit-cognitivi-1}
\begin{center}
  \textbf{Il razionale alla base dell’utilizzo dei dati normativi in caso di interesse al deficit}
\end{center}

Nel caso dell'interesse di valutare il livello funzionale, lo scopo è
valutare il livello della persona: se troppo basso inferiamo che c'è un
deficit.

Nel caso di utilizzo di test neuropsicologici per l'età evolutiva o
nella ricerca di possibili disturbi legati allo sviluppo,
l'intepretazione è stimiamo è se la persona sta avendo una prestazione
che è in linea con il suo sviluppo.

Se è troppo bassa inferiamo che ci sia stato un deficit nello sviluppo.
\end{frame}

\begin{frame}{Valutare i deficit cognitivi}
\phantomsection\label{valutare-i-deficit-cognitivi-2}
\begin{center}
  \textbf{Il razionale alla base dell’utilizzo dei dati normativi}
\end{center}

Esempio:

Ci viene chiesto in contesto clinico di effettuare una valutazione
neuropsicologica di un paziente con trauma cranico.

\emph{Quale potrebbe essere lo scopo della valutazione
neuropsicologica?} \pause

Valutare se è avvenuto un \textbf{cambiamento} (un danno cognitivo), non
valutare le abilità del soggetto in sè.
\end{frame}

\begin{frame}{Utilizzo dei dati normativi per inferire danni cognitivi}
\phantomsection\label{utilizzo-dei-dati-normativi-per-inferire-danni-cognitivi}
\begin{center}
\includegraphics[width=0.15\linewidth,height=\textheight,keepaspectratio]{Figures/triangle.png}
\end{center}

\emph{Nel caso di un ipotetico danno, se avessimo i punteggi ai test
neuropsicologici prima dell'evento lesivo, potremmo confrontarli con
quelli del nostro paziente e vedere se c'è stato un peggioramento
(quindi un danno cognitivo in seguito al trauma).}

Dal momento che questa informazione non è praticamente mai disponibile,
stimiamo il punteggio atteso nel/lla paziente da persone che siano
simili a lui/lei (quindi con i dati normativi).
\end{frame}

\begin{frame}{Valutare i deficit/danni cognitivi}
\phantomsection\label{valutare-i-deficitdanni-cognitivi-2}
\begin{center}
  \textbf{Il razionale alla base dell’utilizzo dei dati normativi}
\end{center}

In particolare: se il punteggio del paziente è molto più basso di quello
che ottengono generalmente persone che sono simili a lui, inferiamo che
c'è un deficit oppure che c'è stato un cambiamento nel tempo cioè che il
paziente ha subito un danno cognitivo.

Per convenzione (e semplicità) definiamo le persone simili a lui sulla
base di età e scolarità (ma è possibile includere anche altre
variabili). In genere questi dati sono già forniti da chi crea il test
sotto forma di \textbf{dati normativi}.
\end{frame}

\begin{frame}{Valutare i deficit cognitivi}
\phantomsection\label{valutare-i-deficit-cognitivi-3}
\begin{center}
  \textbf{Diversi concetti di normalità}
\end{center}

I soggetti che fanno parte dei dati normativi assumiamo che siano
normali, nel senso di autonomi nella vita quotidiana e in assenza di
deficit cognitivi conclamati.

Quando utilizziamo compariamo un punteggio con i dati normativi e
concludiamo che sono nella norma, facciamo un'inferenza statistica a
partire dai punteggi.

Spesso in neuropsicologia si fa riferimento alla curva normale (o
gaussiana) che è una distribuzione con certe proprietà statistiche note
molto utili.

\begin{center}
  \textbf{Sono tre diversi concetti di «normalità»}
\end{center}
\end{frame}

\begin{frame}{Il concetto di normalità nei dati normativi}
\phantomsection\label{il-concetto-di-normalituxe0-nei-dati-normativi}
\small

Da slides precedente ``\emph{I soggetti che fanno parte dei dati
normativi assumiamo che siano \textbf{normali}, nel senso di autonomi
nella vita quotidiana e in assenza di deficit cognitivi conclamati.}''
Ma cosa significa esattamente e in che senso ``normali''?

Teoricamente è riportato con estremo dettaglio cosa si è inteso per
normali. Spesso si intende: autonomi nella vita quotidiana, privi di
patologie conclamate che possono influenzare la cognizione (es. malattie
neurodegenerative o schizofrenia, o disturbi nel neurosviluppo).

In certi casi possono essere più stringenti (es. sopra i 65 anni vengono
richiesti esami medici che escludano segni demenza, es. risonanza,
etc.). A volte anche si riporta testing cognitivo (es. MoCA
\textgreater{} 20). Quest'ultimo però è un aspetto un po' insidioso,
perché i test usati potrebbero a loro volta sbagliare e quindi creare
bias nei dati normativi o circolarità. D'altra parte senza test il
rischio è inserire nel dataset persone che hanno patologie degenerative
in corso.
\end{frame}

\begin{frame}{Tre aspetti fondamentali sui dati normativi 1/3}
\phantomsection\label{tre-aspetti-fondamentali-sui-dati-normativi-13}
\begin{columns}
\column{0.8\textwidth}

\textbf{I dati normativi servono per stimare e classificare la prestazione che ci aspetterebbe da un certo individuo tramite i dati di persone che sono simili all’individuo da valutare.}

\column{0.2\textwidth}

\begin{figure}
\includegraphics[scale=0.05]{Figures/triangle.png}
\end{figure}
\end{columns}
\end{frame}

\begin{frame}{Tre aspetti fondamentali sui dati normativi 2/3}
\phantomsection\label{tre-aspetti-fondamentali-sui-dati-normativi-23}
\begin{columns}
\column{0.8\textwidth}

\textbf{Il fatto che una prestazione sia al di sotto del cut-off di normalità può avere significato o importanza diversa a seconda del contesto e della domanda.}
\vspace{1em}

\small
a) Nel caso di indagine di un disturbo dello sviluppo una performance al di sotto del cut-off è interessante in sé, indicandoci un \textbf{deficit}.

b) nel caso di sospetto di malattia o disturbo acuto nell’adulto viene usata in genere per fare un’inferenza su un peggioramento, quindi un \textbf{danno}, nello stato cognitivo del paziente rispetto ad un evento che può avere una precisa collocazione temporale (es. ictus) oppure no (es. un deterioramento con un decorso legato alla malattia di Alzheimer).
\column{0.2\textwidth}

\begin{figure}
\includegraphics[scale=0.05]{Figures/triangle.png}
\end{figure}
\end{columns}
\end{frame}

\begin{frame}{Tre aspetti fondamentali sui dati normativi 3/3}
\phantomsection\label{tre-aspetti-fondamentali-sui-dati-normativi-33}
\begin{center}
  \textbf{Tre aspetti fondamentali su cut-off e dati normativi (3)}
\end{center}
\vspace{1em}

\begin{columns}
\column{0.8\textwidth}
\small

\textbf{In Neuropsicologia forense, l’interesse nel deficit o nel danno, dipende dal contesto e dalla domanda:}
\vspace{1em}

\small
a) nel caso di una perizia per valutare un danno in seguito ad un evento (es. perizia per valutare disturbi cognitivi in seguito ad un trauma cranico), il confronto con dati normativi è utili per inferire un danno (inteso come peggioramento rispetto a condizione pre-esistente)

b) in caso di valutazioni per valutare il lvello del paziente indipendentemente dalla causa (es. valutare capacità di intendere e volere) allora è interessante usare i dati normativi per una classificazione della performance in sé.
\column{0.2\textwidth}

\begin{figure}
\includegraphics[scale=0.05]{Figures/triangle.png}
\end{figure}
\end{columns}
\end{frame}

\begin{frame}{Raccolta di dati normativi}
\phantomsection\label{raccolta-di-dati-normativi}
\begin{center}
\begin{tikzpicture}

% --- Your figure -------------------------------------------------------
\node (img) at (0,0) {\includegraphics[width=0.7\textwidth]{Figures/dati_normativi.png}};

% Top-left
\only<1>{\fill[white] (-3.5,3) rectangle (0,0);}

% Top-right
\only<1-2>{\fill[white] (-0.2,3) rectangle (3,0);}

% Bottom-left
\only<1-3>{\fill[white] (-3.5,0) rectangle (0,-3);}

% Bottom-right
\only<1-4>{\fill[white] (-0.2,0) rectangle (3,-3);}


\end{tikzpicture}
\end{center}
\onslide<5> % <-- forces slide 5 to exist
La linea verticale indica il \textbf{cut-off di normalità}.

\pause
\end{frame}

\begin{frame}{Il significato di cut-off di normalità}
\phantomsection\label{il-significato-di-cut-off-di-normalituxe0}
\begin{columns}
\column{0.8\textwidth}

Il cut-off è un punteggio particolarmente basso che osserviamo in pochi soggetti normali e che pertanto inferiamo sia segno di un deficit/danno cognitivo.
\vspace{1em}

In genere il cut-off è definito (arbitrariamente) come il valore che delimita il 5\% dei punteggi più bassi.
\vspace{1em}

NOTA: Il cut-off NON ci indica la probabilità che il soggetto sia normale o abbia un danno cognitivo (vedi slides Valutare i deficit cognitivi e vedi Bayes, è una differenza un po’ sottile)

\column{0.2\textwidth}

\begin{figure}
\includegraphics[scale=0.05]{Figures/triangle.png}
\end{figure}
\end{columns}
\end{frame}

\begin{frame}{Un fondamentale aspetto sull'utilizzo dei dati normativi.}
\phantomsection\label{un-fondamentale-aspetto-sullutilizzo-dei-dati-normativi.}
\begin{columns}
\column{0.8\textwidth}

\small
Nella mia definizione di dati normativi un elemento importante è che il confronto deve essere effettuate con persone simili alla persona che devo testare.
\vspace{1em}

Ma come definire se sono simili?
\vspace{1em}

Come potremmo definire la popolazione adeguata per un confronto normativo di una persona con 68 anni e ha fatto le scuole elementari e lavora come contadino?
Gli italiani con età tra 65 e 70 anni?
Gli italiani che hanno esattamente 68 anni e chehanno fatto solo le scuole elementari? Che lavorano come contadini?
\vspace{1em}

\textbf{Non esiste un criterio univoco per definire cosa significhi essere simili, ma più sono gli elementi di somiglianza migliore sarà il nostro confronto}

\column{0.2\textwidth}

\begin{figure}
\includegraphics[scale=0.05]{Figures/triangle.png}
\end{figure}
\end{columns}
\end{frame}

\begin{frame}{Valutare i deficit cognitivi}
\phantomsection\label{valutare-i-deficit-cognitivi-4}
\begin{center}
  \textbf{Un altro aspetto fondamentale su cut-off e dati normativi}
\end{center}
\vspace{2em}

\begin{columns}
\column{0.8\textwidth}

A livello di operazioni da svolgere, la stima un deficit vs la stima di un danno sono simili a livello statistico: entrambe implicano il confrontare la prestazione con i dati normativi (vedi slides precedente per scelta appropriata del confronto).

Quello che cambia è l’\textbf{interpretazione} (vedi slides su Interpretazione) che in un caso può essere quella di un deficit, in un caso quella di un danno.

\column{0.2\textwidth}

\begin{figure}
\includegraphics[scale=0.05]{Figures/triangle.png}
\end{figure}
\end{columns}
\end{frame}

\begin{frame}{Una riflessione su confronto con dati normativi}
\phantomsection\label{una-riflessione-su-confronto-con-dati-normativi}
\begin{columns}
\column{0.8\textwidth}

Questo è molto importante a livello forense, perché la performance del test confrontata con i cut-off in sé non ci dice se il paziente è peggiorato prima dell’evento rilevante (es. il paziente ha deficit per il TBI) oppure se è un deficit già presente (è un disturbo di neurosviluppo). 

\textbf{Il confronto col cut-off ci dice osservare quella prestazione è molto improbabile, se il paziente fosse sano (come le persone dei dati normativi).}

\column{0.2\textwidth}

\begin{figure}
\includegraphics[scale=0.05]{Figures/triangle.png}
\end{figure}
\end{columns}
\end{frame}

\begin{frame}{Uno o più cut-off? 1/2}
\phantomsection\label{uno-o-piuxf9-cut-off-12}
\small

Una volta che si possiede una distribuzione di dati di riferimento è
possibile in generale classificare la presentazione rispetto all'intera
distribuzione (questo può dipendere dal metodo con cui sono analizzati i
dati).

Ad esempio è possibile dire se la persona ha una performance più bassa
del 20\% o il 30\% rispetto al campione/popolazione di riferimento
(approfondimento nella sezione: Valutare il Deficit cognitivo Formule).

L'assenza di soglie ci porta però (da un punto di vista concettuale) ad
un rischio del
\href{https://it.wikipedia.org/wiki/Paradosso_del_sorite}{\ul{paradosso
del mucchio}}.

Non è possibile definire chiaramente quando si passa da normalità a
quando si passa alla presenza di un ``deficit/danno''. Il paradosso del
mucchio ci insegna che \textbf{le categorie a partire da una variabile
continua sono e saranno sempre arbitrarie}.
\end{frame}

\begin{frame}{Uno o più cut-off? 2/2}
\phantomsection\label{uno-o-piuxf9-cut-off-22}
Il fatto di avere una soglia decisionale, per quanto arbitraria, ha però
delle utilità euristiche. Dovremo decidere ad un certo punto se c'è un
problema oppure no e questo tendenzialmente ci porta ad utilizzare
soglie categoriali.

Si veda per esempio
\href{https://journals.sagepub.com/doi/full/10.1177/09593543231160112}{\underline{questo articolo}}
per un approfondimento sul valore epistemico di scelte dicotomiche, che
ben si applica a questo caso.
\end{frame}

\begin{frame}{Perché questo modo ``contorto'' di definire il deficit
cognitivo e il danno?}
\phantomsection\label{perchuxe9-questo-modo-contorto-di-definire-il-deficit-cognitivo-e-il-danno}
Il modo in cui definiamo danno e deficit (a partire da soggetti normali)
potrebbe sembrare contorto. Perchè dobbiamo basarci su questo metodo?

Se avessimo un criterio oggettivo (es. danno cerebrale) su cui basarci,
sarebbe concettualmente più semplice. Potremmo usare quello e usare i
test neuropsicologici come strumenti per identificare questo tipo di
danno (entreremmo nel reame della validità di criterio, sensibilità,
specificità, etc. vedi slides su Validità e slides su Sensibilità e
Specificità).
\end{frame}

\begin{frame}{Perché questo modo ``contorto'' di definire il deficit
cognitivo e il danno?}
\phantomsection\label{perchuxe9-questo-modo-contorto-di-definire-il-deficit-cognitivo-e-il-danno-1}
Il problema è metodologico/epistemologico. Come possiamo definire un
deficit?

Come possiamo essere certi che ci sia un deficit cognitivo?

Il problema è che non è possibile definire in maniera non circolare il
deficit cognitivo e per tale ragione usiamo questo metodo \emph{alla
rovescia}. Definiamo cosa è l'\textbf{assenza di deficit cognitivo} (su
base normativa), e inferiamo la presenza di deficit a partire da questo.
\end{frame}

\begin{frame}{Perché questo modo ``contorto'' di definire il deficit
cognitivo e il danno?}
\phantomsection\label{perchuxe9-questo-modo-contorto-di-definire-il-deficit-cognitivo-e-il-danno-2}
Il problema è in questo caso di assenza di dati precedenti che mi
possano aiutare a capire se c'è stato un peggioramento o un effettivo
danno cognitivo.

Usare un dato osservabile (es. danno cerebrale) sarebbe rischioso e
limitante. In molti casi può esserci danno cognitivo senza apparente
danno cerebrale o viceversa (es. nel trauma cranico).
\end{frame}

\begin{frame}{Perché questo modo ``contorto'' di definire il deficit
cognitivo e il danno?}
\phantomsection\label{perchuxe9-questo-modo-contorto-di-definire-il-deficit-cognitivo-e-il-danno-3}
In tal caso stimiamo la presenza di un danno se la persona ha un
punteggio anomalo rispetto al campione normativo. In questo caso il
passaggio è doppio:

\begin{enumerate}
[a)]
\item
  Sto assumendo che la persona è ben rappresentata dal campione
  normativo
\item
  sto inferendo che se la sue performance è bassa, probabilmente c'è
  stato un danno/peggioramento.
\end{enumerate}
\end{frame}

\begin{frame}{Perché questo modo ``contorto'' di definire il deficit
cognitivo e il danno?}
\phantomsection\label{perchuxe9-questo-modo-contorto-di-definire-il-deficit-cognitivo-e-il-danno-4}
Il fatto di definire un \textbf{danno} (quindi un peggioramento) a
partire dal confronto con i dati normativi è una strategia per
compensare il fatto che non abbiamo dati precedenti al sospetto evento
che potrebbe avere causato il peggioramento.

Ma attenzione a ricordare che l'uso è solo una strategia con i suoi
limiti, basata su una stima attesa della performance del paziente.
\end{frame}

\begin{frame}{Una nota sul definire il danno: Esempio 1}
\phantomsection\label{una-nota-sul-definire-il-danno-esempio-1}
Supponiamo di avere un test in cui il cut-off è 10 (5\%) con elevata
affidabilità e nessun effetto pratica. Supponiamo di testare un paziente
con trauma cranico per valutare il danno.

Il paziente ha ottenuto 14 (sopra il cut-off).

Potrebbe essere che prima dell'evento aveva un punteggio di 20 e il
punteggio di 14 potrebbe riflettere un cambiamento non colto dal
confronto con il cut-off di 10 (vedi anche slides valutare il
cambiamento nel tempo).
\end{frame}

\begin{frame}{Una nota sul definire il danno: Esempio 2}
\phantomsection\label{una-nota-sul-definire-il-danno-esempio-2}
Supponiamo di avere un test in cui il cut-off è 10 (5\%) con elevata
affidabilità e nessun effetto pratica. Supponiamo di testare un paziente
con trauma cranico per valutare il danno.

Il paziente ha ottenuto 8 (sotto il cut-off).

Il paziente aveva 8 anche prima del trauma cranico, in realtà non è
peggiorato (vedi anche slides valutare il cambiamento nel tempo).
\end{frame}

\begin{frame}{Una nota sul definire il danno}
\phantomsection\label{una-nota-sul-definire-il-danno}
\small

Il cut-off non era adeguato perché il nostro paziente testato era un
paziente con prestazione presumibilmente già estrema (molto alta o
bassa) rispetto ai dati normativi.

Per la \textbf{valutazione del danno} i dati normativi funzionano bene
(per come sono disegnati) se la persona è ben rappresentata da essi e se
non aveva (prima dell'ipotetico evento legato al danno), già una
prestazione agli estremi (bassa o alta).

\emph{In assenza di ogni altra informazione}, questa è un'assunzione
ragionevole.

Assumiamo che la nostra persona abbia avuto (prima del danno) una
prestazione analoga alla maggior parte delle persone simili a lui. Altre
informazioni (es. da anamnesi, da aspetti qualitativi, da storia
personale), potrebbero suggerire che invece la persona non è ben
rappresentata.
\end{frame}

\begin{frame}{Valutare i deficit cognitivi}
\phantomsection\label{valutare-i-deficit-cognitivi-5}
\begin{center}
  \textbf{Una nota sul definire il danno}
\end{center}

Il fatto di avere un punteggio sotto (o sopra) cut-off è solo
un'informazione legata ad alcune \emph{probabilità} (poche persone
simili al paziente testato ottengono questa performance), ma le
conclusioni che ne traiamo sono \emph{interpretazioni}.

Altre informazioni a disposizione potrebbero farci propendere per una o
l'altra intepretazione, se sono sufficientemente forti per farci
ritenere che il confronto con i dati normativi non è adeguato o dovrebbe
essere rivisto.

Ovviamente questi ragionamento non delegittima il confronto i cut-off
(che rimane uno strumento solido), ma ne sottolineano alcuni potenziali
limiti.
\end{frame}

\begin{frame}{Valutare i deficit cognitivi}
\phantomsection\label{valutare-i-deficit-cognitivi-6}
\begin{center}
  \textbf{Uno o più cut-off?}
\end{center}

Per questo motivo e per motivi di praticità spesso il fatto di
concludere che c'è un danno è legato al cut-off di normalità citato in
precedenza.

Per un neuropsicologo clinico, laddove sia possibile a partire dai
metodi statistici a disposizione è comunque utile sapere la
classificazione esatta (es. 20\%, 14\% della popolazione di riferimento
e non semplicemenete sopra/sotto cut-off) perchè potrebbe essere
comunque utile e informativa.

Es. una performance che delimita è al 5.1\% del campione/popolazione di
riferimento, sarebbe comunque sopra cut-off (esattamente come una al
95\%).
\end{frame}

\begin{frame}{Come si calcolano i cut-offs?}
\phantomsection\label{come-si-calcolano-i-cut-offs}
\begin{itemize}
\tightlist
\item
  Metodi parametrici (Media, deviazione standard e Z score)
\item
  Metodi non parametrici
\item
  Altri metodi (regressioni)
\end{itemize}
\end{frame}

\begin{frame}{Valutare i deficit cognitivi}
\phantomsection\label{valutare-i-deficit-cognitivi-7}
\begin{center}
  \textbf{Metodi parametrici (base) per il calcolo dei cut-off}
\end{center}

\begin{enumerate}
[1)]
\item
  \textbf{Media e deviazione standard}. Se la distribuzione è normale
  possiamo calcolare il cut-off del 5\% dei punteggi più bassi, come
  1.64 deviazioni standard sotto la media.
\item
  \textbf{Z-score}. Metodo analogo al precendete. Se trasformiamo i dati
  in punti \(Z\), un punteggio pari a -1.64 indica il cut-off dei 5\%
  dei punteggi più bassi.
\end{enumerate}

\[
Z = \frac{X - M}{SD}
\] \[
\small
\begin{aligned}
X & = \text{punteggio ottenuto dal soggetto} \\
M & = \text{media del gruppo di riferimento} \\
SD & = \text{deviazione standard del gruppo di riferimento} \\
\end{aligned}
\]
\end{frame}

\begin{frame}{Valutare i deficit cognitivi}
\phantomsection\label{valutare-i-deficit-cognitivi-8}
\begin{center}
  \textbf{Metodi parametrici (base) per il calcolo dei cut-off}
\end{center}
\vspace{1em}

\begin{center}
\includegraphics[width=0.8\linewidth,height=\textheight,keepaspectratio]{Figures/gaussiana.png}
\end{center}
\small \emph{NOTA}: riferirci alla curva normale e alla trasformazione
in punti \(Z\) ci aiuta ad identificare percentili, cioè valori che
delimitano percentuali di dati.

\emph{I metodi parametrici ossono essere usati con sicurezza solo se la
distribuzione è normale (gaussiana)}
\end{frame}

\begin{frame}{Valutare i deficit cognitivi}
\phantomsection\label{valutare-i-deficit-cognitivi-9}
\begin{columns}
\column{0.5\textwidth}

\begin{center}
\textbf{distribuzione normale}
\end{center}
\vspace{1em}
\begin{figure}
\includegraphics[scale=0.25]{Figures/distr_norm.png}
\end{figure}
\textbf{ }

\column{0.5\textwidth}

\begin{center}
\textbf{distribuzione non-normale}
\end{center}
\vspace{1em}
\begin{figure}
\includegraphics[scale=0.25]{Figures/distr_nonnorm.png}
\end{figure}

\begin{center}
test con \textbf{effetto soffitto}
\end{center}

\end{columns}
\end{frame}

\begin{frame}{Valutare i deficit cognitivi}
\phantomsection\label{valutare-i-deficit-cognitivi-10}
\begin{columns}
\column{0.8\textwidth}
Se si usano i punti z con distribuzioni non normali, la corrispondenza tra specifici punteggi z e percentuali può non essere rispettata.
\vspace{2em}

Ad esempio -1.64 può non corrispondere al valore che delimita 5\% dei punteggi più bassi, portando al calcolo di cut-off sbagliati.
\column{0.2\textwidth}

\begin{figure}
\includegraphics[scale=0.05]{Figures/triangle.png}
\end{figure}
\end{columns}
\end{frame}

\begin{frame}{Valutare i deficit cognitivi}
\phantomsection\label{valutare-i-deficit-cognitivi-11}
\begin{center}
  \textbf{PARENTESI}
  
  \textbf{Standardizzare (trasformare in z-score) a volte è anche detto "normalizzare", ma non rende la distribuzione \textit{normale}!}
\end{center}

\begin{center}
\includegraphics[width=1\linewidth,height=\textheight,keepaspectratio]{Figures/distribuzioni_normalizzate.png}
\end{center}
\end{frame}

\begin{frame}{Valutare i deficit cognitivi}
\phantomsection\label{valutare-i-deficit-cognitivi-12}
\begin{center}
  \textbf{Metodi non parametrici (base) per il calcolo dei cut-off}
\end{center}
\vspace{2em}

I \textbf{percentili non parametrici} sono il metodo più diffuso di
calcolo dei cut-off nel caso di distribuzioni non normali.

Lo svantaggio dell'utilizzo di percentili è che non è possibile usare
formule rapide per ricavare cut-off a diverse percentuli (ad esempio il
20\% dei punteggi peggiori) ma è necessario avere a disposizione tutti i
dati.
\end{frame}

\begin{frame}{Valutare i deficit cognitivi}
\phantomsection\label{valutare-i-deficit-cognitivi-13}
\begin{center}
  \textbf{Altri metodi (avanzati) per il calcolo dei cut-off}
\end{center}
\vspace{1em}

Esistono altri metodi per il calcolo dei cut-off. La maggior parte
prevedono l'uso di regressioni.

Il vantaggio, rispetto a metodi che prevedono la divisione in fasce è
che con questi metodi \textbf{non avviene una stratificazione
arbitraria}.

Esempi:

\begin{itemize}
\item
  metodo dei Punteggi Equivalenti (Capitani \& Laiacona, 2007)
\item
  Alcuni metodi di Crawford (Crawford \& Howell, 1998; Crawford \&
  Garthwaite, 2006)
\end{itemize}
\end{frame}

\begin{frame}{Campione e Popolazione nei dati normativi}
\phantomsection\label{campione-e-popolazione-nei-dati-normativi}
I dati normativi sono i dati raccolti da un \textbf{campione}
proveniente da una \textbf{popolazione}.

L'interesse è però sempre confrontare la prestazione osservata con i
dati della popolazione. Alcuni metodi tengono conto di questa
differenza, fornendo dei cut-off che sono una stima di quelli che si
osserverebbero nella popolazione.

\textbf{Questi cut-off sono spesso più bassi} (meno restrittivi).

Il metodo dei \textbf{Punteggi Equivalenti} (Capitani e Laiacona, 2007)
e il metodo di \textbf{Crawford} (Crawford /\& Howell, 1998; Crawford
/\& Garthwaite, 2006) tengono conto di questa differenza. Tali metodi
sono da considerarsi più corretti.
\end{frame}

\begin{frame}{Variabili demografiche nei dati normativi}
\phantomsection\label{variabili-demografiche-nei-dati-normativi}
I dati normativi spesso calcolati separatamente per età e scolarità o
altre variabili demografiche, per poter confrontare adeguatamente la
prestazione di un determinato individuo.

I metodi che utilizzano regressioni cercano di modellare l'effetto di
età e scolarità e tenerne conto.
\end{frame}

\begin{frame}{Variabili demografiche nei dati normativi}
\phantomsection\label{variabili-demografiche-nei-dati-normativi-1}
I metodi «base» descritti di solito dividono il campione in vari
sotto-campioni e calcolano diversi cut-off per ciascun sotto-campione.

\begin{itemize}
\item
  In genere sono forniti valori separati per diverse combinazioni di età
  e scolarità, suddivise in fasce.
\item
  Possono essere fornite (ma raramente) tabelle di correzioni, o a volte
  singole correzioni (ad esempio +2 per scolarità \textless= 8).
\end{itemize}
\end{frame}

\begin{frame}{Variabili demografiche nei dati normativi}
\phantomsection\label{variabili-demografiche-nei-dati-normativi-2}
I metodi che utilizzano regressioni non suddividono il campione totale.
In genere forniscono i cut-off in diverse modalità.

\begin{itemize}
\item
  Un singolo cut-off e tabelle di correzione.
\item
  Un singolo cut-off e formule incomprensibili per calcolare la
  correzione.
\item
  Tabelle con diversi cut-off, per le varie combinazioni di età,
  scolarità o sesso.
\end{itemize}
\end{frame}

\begin{frame}{Variabili demografiche nei dati normativi}
\phantomsection\label{variabili-demografiche-nei-dati-normativi-3}
\begin{center}
\includegraphics[width=0.8\linewidth,height=\textheight,keepaspectratio]{Figures/cutoff_humor.png}
\end{center}

\[
\scalebox{0.85}{$
sd_{Y_i-\hat{Y_i}} = \frac{SE_{Y-\hat{Y}}}{\sqrt{n}}\sqrt{n+1+\frac{z_{i,age}^2}{1-R_{age}^2}+\frac{z_{i,education}^2}{1-R_{i,education}^2}-2\frac{\beta_{education,age}z_{i,age}z_{i,education}}{1-R_{education}^2}}
$}
\]
\end{frame}

\begin{frame}{Quale è il metodo migliore per il calcolo dei cut-off?}
\phantomsection\label{quale-uxe8-il-metodo-migliore-per-il-calcolo-dei-cut-off}
\pause

Dipende.

\pause

Metodi che utilizzano distinizione campione/popolazione sono
preferibili.

Attenzione all'uso indiscriminato dei punti-z (se non è nota la
distribuzione punteggi)

Attenzione ai test che usano metodi basati su regressioni. Talvolta
possono portare a stime distorte, specie per valori delle variabili (età
e scolarità) che sono estremi (es. bassa scolarità o alta età).
\end{frame}

\begin{frame}{Riepilogo su dati normativi e cut-off}
\phantomsection\label{riepilogo-su-dati-normativi-e-cut-off}
\small

\begin{itemize}
\item
  I dati normativi servono per stimare la prestazione che ci
  aspetterebbe da un certo individuo tramite i dati di persone che sono
  simili all'individuo da valutare.
\item
  Il confronto con dati normativi è fatto per effettuare un'inferenza su
  un possibile deficit cognitvo (disturbo funzioanle, che sia legato a
  danno o sviluppo o pre- esistente) o danno (peggioramento) dello stato
  cognitivo del paziente. Spesso i metodi usati sono uguali, quello che
  cambia e l'interpretazione
\item
  Il cut-off (in senso generale) è definito come il valore che delimita
  il 5\% dei punteggi più bassi.
\item
  Esistono vari modi di calcolare i cut-off
\end{itemize}
\end{frame}

\begin{frame}{Aspetti da considerare nell'utilizzo dei dati normativi}
\phantomsection\label{aspetti-da-considerare-nellutilizzo-dei-dati-normativi}
Come principio generale, per stimare correttamente la prestazione del
mio soggetto (e rendere sensato il confronto con i dati normativi), i
dati normativi devono essere rappresentativi. Di seguito sono esaminate
alcune comuni domande.
\end{frame}

\begin{frame}{Aspetti da considerare nell'utilizzo dei dati normativi
(1)}
\phantomsection\label{aspetti-da-considerare-nellutilizzo-dei-dati-normativi-1}
\textbf{1) Utilizzare dati normativi aggiornati} \vspace{2em}

Dati normativi raccolti molti anni fa possono non essere adeguati.
Questo è legato al cosiddetto Flynn effect (Flynn, 1987; Williams, 2013)
\end{frame}

\begin{frame}{Aspetti da considerare nell'utilizzo dei dati normativi
(2)}
\phantomsection\label{aspetti-da-considerare-nellutilizzo-dei-dati-normativi-2}
\textbf{2) Specie per stimare il danno è importante che gli individui
dei campioni normativi siano simili all'individuo da valutare. (campione
ampio, ma non simile non è utile).} \vspace{2em}

Un campione molto numeroso ma poco rappresentativo non è utile.
\end{frame}

\begin{frame}{Aspetti da considerare nell'utilizzo dei dati normativi}
\phantomsection\label{aspetti-da-considerare-nellutilizzo-dei-dati-normativi-3}
\begin{columns}
\column{0.8\textwidth}
\small
\textbf{3) Utilizzare dati normativi paese-specifici.}
\vspace{1em}

Non si può stimare la prestazione a partire dai dati di soggetti che provengono da altri paesi. \textbf{Non basta tradurre un test per renderlo utilizzabile.}

Non è solo un fattore di lingua e di adattamento (che già minano la validità del test) ma anche legato ad aspetti socioculturali che possono rendere i dati normativi non adeguati. Non basta avere stessa età e la scolarità a rendere un soggetto comparabile con dati normativi.
L’iniziale diffusione del \textbf{MoCA} (Nasreddine et al., 2005) in Italia, è un caso lampante di questa prassi profondamente scorretta.

\column{0.2\textwidth}

\begin{figure}
\includegraphics[scale=0.05]{Figures/triangle.png}
\end{figure}
\end{columns}
\end{frame}

\begin{frame}{Aspetti da considerare nell'utilizzo dei dati normativi}
\phantomsection\label{aspetti-da-considerare-nellutilizzo-dei-dati-normativi-4}
\begin{columns}
\column{0.8\textwidth}
\small
\textbf{3) Utilizzare dati normativi paese-specifici.}
\vspace{1em}

Il razionale di uso dati normativi è quello di confrontare la prestazione con prestazione di persone “simili alla persona che intendiamo valutare”. 

Usare norme da altri paesi è già in sé una grande differenza. Confrontando Norme esistenti tra paesi questo può essere verificato empiricamente. Spesso hanno valori di soglia (e cut-offs) molto diversi.

\column{0.2\textwidth}

\begin{figure}
\includegraphics[scale=0.05]{Figures/triangle.png}
\end{figure}
\end{columns}
\end{frame}

\begin{frame}{Riassunto}
\phantomsection\label{riassunto}
In neuropsicologia clinica spesso siamo chiamati valutare due aspetti:

C'è un \textbf{deficit cognitivo} o dobbiamo stimare il livello di
abilità cognitiva generale del paziente? \emph{Questo è fatto
confrontando con dati normativi e classificando (spesso con percentuali
o sopra/sotto norma)}

C'è stato un \textbf{danno cognitivo}?, cioè un peggioramento? Anche
questo viene fatto con i dati normativi, ma con una serie di assunzioni
aggiuntive nell'\underline{interpretazione}.

Ci sono vari metodi matematici (con pro e contro e più o meno robusti),
per definire i cut-off ed eventuali altre soglie.
\end{frame}




\end{document}
